% !TeX root = ../main.tex

\setproblem{G}{\ProblemTitleG}
\section{\texorpdfstring{\pletter. \ptitle}{문제 \pletter}}

\begin{frame}[label=problem-G]
    \sectiontitle{\pletter}{\ptitle}
    \sectionmeta{
        출제자: \idtext{areian13}\\
        태그: \tagtext{\#bruteforcing, \#implementation}\\
        예상 난이도: \tiertext{Bronze2}
    }
    \begin{itemize}
        \item 제출 20번, 정답 11명 (정답률 55\%)
        \item 처음 푼 사람: ㅇㅎㅈ (00:16:47)
    \end{itemize}
\end{frame}

\begin{frame}{\textbf{\pletter}. \ptitle}
    \textbf{사진 두 장을 직접 골라 보기}
\begin{itemize}
        \item 사진이 최대 $100$장이므로 어떤 사진을 고를지 모두 확인할 수 있습니다.
        \item 사진 $i,j$를 골랐다면, 각 부원이 두 사진 중 적어도 하나에서 눈을 떴는지 셉니다.
        \item 문자열의 같은 위치를 비교해 둘 중 하나라도 \texttt{1}이면 $1$명으로 셉니다.
        \item 두 사진 모두 눈을 떴어도 같은 사람이므로 한 번만 세어야 합니다.
    \end{itemize}
    \begin{center}
        \begin{tabular}{c|c}
            첫 번째 사진 & \texttt{10100}\\
            두 번째 사진 & \texttt{01100}\\\hline
            적어도 한 장에서 눈을 뜬 부원 & \texttt{11100}: $3$명
        \end{tabular}
    \end{center}
\end{frame}

\begin{frame}{\textbf{\pletter}. \ptitle}
    \textbf{한 장만 고르는 경우도 포함하기}
\begin{itemize}
        \item $i$를 $1$부터 $N$까지, $j$를 $i$부터 $N$까지 바꾸며 각 선택의 점수를 구합니다.
        \item $i=j$인 경우는 사진 한 장만 고른 것으로 보면 됩니다.
        \item 이렇게 하면 사진이 한 장뿐인 경우도 따로 처리할 필요가 없습니다.
        \item 각 선택에서 얻은 점수 중 가장 큰 값을 출력합니다.
        \item 사진 쌍이 $O(N^2)$개이고 부원 $M$명을 확인하므로 시간 복잡도는 $O(N^2M)$입니다.
        \item 최대 약 $10^6$번의 문자 비교이므로 충분합니다.
    \end{itemize}
\end{frame}
